\section{Modelo dinámico}

\subsection{Sistema y fuentes de CO$_2$}

El sistema experimental se representa mediante dos fuentes acopladas: la respiración larvaria y la degradación microbiana del alimento o lecho. La tasa total de producción de CO$_2$ es
\begin{equation}
  R_{\mathrm{total}}(t)
  =
  R_{\mathrm{lar}}(t)+R_{\mathrm{mic}}(t).
  \label{eq:fuentes-totales}
\end{equation}
Esta descomposición permite contrastar por separado los controles sin larvas, las larvas en ayuno y los tratamientos completos antes de evaluar el sistema integrado.

\subsection{Estados larvarios}

La larva se describe mediante tres compartimentos: asimilados disponibles $A(t)$, biomasa estructural $B(t)$ y reservas lipídicas $L(t)$. El crecimiento medido corresponde a una observación agregada de la biomasa larvaria, de modo que no identifica por sí solo cada compartimento.

\subsection{Alimentación y asimilación}

La ingesta dependerá de la dieta, la disponibilidad del sustrato y el estado de desarrollo. La fracción asimilada alimenta el compartimento $A$, desde el cual se atienden mantenimiento, síntesis de biomasa estructural y acumulación de lípidos. La formulación completa se fijará después de verificar las unidades de todos los parámetros.

\subsection{Crecimiento y respiración larvaria}

La producción larvaria de CO$_2$ se descompone en mantenimiento y costos respiratorios de síntesis:
\begin{equation}
  R_{\mathrm{lar}}(t)
  =
  r_{\mathrm{m}}(t)
  +\gamma_B\,\dot B(t)
  +\gamma_L\,\max\{\dot L(t),0\}.
  \label{eq:respiracion-larvaria}
\end{equation}
No se asigna producción de CO$_2$ a la movilización neta de lípidos; esa contribución queda asociada con el mantenimiento y con los flujos de síntesis efectivamente positivos.

\subsection{Módulo microbiano del lecho}

El lecho se representa mediante materia orgánica degradable $D_M(t)$ y una variable microbiana efectiva $X(t)$. La degradación sigue una cinética de disponibilidad:
\begin{equation}
  r_{\mathrm{deg}}(t)
  =
  k_{\mathrm{deg}}X(t)
  \frac{D_M(t)}{K_{SM}+D_M(t)},
  \qquad
  R_{\mathrm{mic}}(t)
  =
  Y_{\mathrm{CO_2}}r_{\mathrm{deg}}(t).
  \label{eq:modulo-microbiano}
\end{equation}
Los controles de alimento sin larvas son la restricción principal de este módulo.

\subsection{Balance de la cámara de medición}

La variable observada por el sensor es la pendiente temporal de la concentración en la cámara cerrada. Si $R_{\mathrm{total}}$ está en mol por unidad de tiempo,
\begin{equation}
  \frac{dC_{\mathrm{ppm}}}{dt}
  =
  10^6\,\frac{R_uT}{P_{\mathrm{atm}}V_{\mathrm{air}}}\,
  R_{\mathrm{total}}(t).
  \label{eq:balance-camara}
\end{equation}
Las tasas experimentales se recalcularán intervalo por intervalo a partir de las marcas temporales reales; no se aplicará un factor global de conversión.

\subsection{Observables e identificabilidad}

\begin{center}
\begin{tabular}{lll}
\toprule
Medición & Fuente dominante & Restricción del modelo \\
\midrule
Alimento D1 sin larvas & $R_{\mathrm{mic}}$ & Módulo microbiano D1 \\
Alimento D4 sin larvas & $R_{\mathrm{mic}}$ & Módulo microbiano D4 \\
Larvas solas en ayuno & $R_{\mathrm{lar}}$ & Mantenimiento larvario \\
Alimento D1 con larvas & $R_{\mathrm{total}}$ & Acoplamiento en D1 \\
Alimento D4 con larvas & $R_{\mathrm{total}}$ & Acoplamiento en D4 \\
Crecimiento D1 y D4 & Biomasa larvaria & Crecimiento agregado \\
\bottomrule
\end{tabular}
\end{center}

La diferencia tratamiento--control no se interpretará automáticamente como respiración larvaria, porque las larvas también pueden modificar la actividad microbiana del lecho. La validación debe combinar componentes independientes, sistema integrado y transferencia a datos externos.
